\documentclass[journal, letterpaper]{IEEEtran}
\usepackage{lipsum}
\usepackage{titlesec}
\usepackage[table]{xcolor}
\usepackage{makecell}
\usepackage[most]{tcolorbox}


\titleformat{\section}
    {\centering\normalfont\Large\bfseries}
    {}
    {}
    {}

\newtcolorbox{theory}[2][]{breakable,sharp corners, skin=enhancedmiddle jigsaw,
parbox=false, boxrule=0mm, leftrule=2mm, boxsep=0mm, arc=0mm, outer arc=0mm,
attach title to upper, after title={\ },
coltitle=black, coltext=black,
colback=blue!10, colframe=blue!50,
title={#2}, fonttitle=\bfseries, #1}

\newtcolorbox{example}[2][]{breakable,sharp corners, skin=enhancedmiddle jigsaw,
parbox=false, boxrule=0mm, leftrule=2mm, boxsep=0mm, arc=0mm, outer arc=0mm,
attach title to upper, after title={\ },
coltitle=black, coltext=black,
colback=gray!10, colframe=gray!50,
title={#2}, fonttitle=\bfseries, #1}

\newtcolorbox{aside}[2][]{breakable,sharp corners, skin=enhancedmiddle jigsaw,
parbox=false, boxrule=0mm, leftrule=2mm, boxsep=0mm, arc=0mm, outer arc=0mm,
attach title to upper, after title={\ },
coltitle=black, coltext=black,
colback=red!10, colframe=red!50,
title={#2}, fonttitle=\bfseries, #1}


\title{COMP3121 Notes}
\author{Ronnie Low}
\date{Term 1, 2026}

\begin{document}

\maketitle

\begin{theory}{What is an algorithm?} \\
    Algorithms are a set of precisely defined steps are executed based off mechanical steps - there is no intelligence involved, just follow the steps given.
\end{theory}

\begin{aside}{Types of algorithms} \\
    \textit{1. Sequential deterministic algorithms:}
        \begin{itemize}
            \item Only one step can be executed at a time
            \item No element of randomness
        \end{itemize}
    \textit{2. Parallel algorithms:}
        \begin{itemize}
            \item Multiple steps can be ran at a time
            \item Probably not covered in 3121
        \end{itemize}
    \textit{3. Randomised algorithms:}
        \begin{itemize}
            \item Utilises random values
            \item Used in quantum computing
            \item Probably not covered in 3121
        \end{itemize}
\end{aside}

\begin{theory}{Why prove algorithms?} \\
    Sometimes, its not obvious if an algorithm does the things its advertised through common sense:
    \begin{itemize}
        \item Will your algorithm terminate in a timely manner?
        \item Can you prove that the algorithm won't infinitely loop?
        \item Is the time complexity as advertised?
    \end{itemize}
    Proofs are needed for these cases, most of the time they're the \textbf{only} way to know if the algorithm does its job.
\end{theory}

\begin{example}{Stable matching example} \\
    There are \textit{n} hospitals and \textit{n} doctors. Each hospital wants to hire exactly one doctor. Each doctor and hospital submit a list of preferences ranking which hospital they want to go to and which doctor they want.

    How do we design an algorithm where there is no hospital doctor pair (\textit{h, d}) that are dissatisfied with their result (stable):

    \begin{center}
        for pair $p = (h, d')$ \& $p' = (h', d)$\\
        \begin{itemize}
            \item $h$ prefers $d$ over $d'$
            \item $h'$ prefers $d'$ ovr $d$
        \end{itemize}
    \end{center}
    This implies that both hospitals and doctors have a preferred swap to make (unstable
    match) \par
    \bigskip
    \noindent\textbf{Stable matching algorithm} \par
    \noindent While there is a free hospital, pick a free hospital $h$ and offer a job to the highest
    ranking doctor $d$ on its preference list. \par
    \bigskip
    \noindent If $d$ doesn't have a job yet, $d$ always accepts the job.\par
    \par
    else if $d$ is already in a pair $(h' d)$:
    \begin{itemize}
        \item{if $h$ is higher ranked on the preference list for $d$, a new pair $(h, d)$ is
            created and the previous pair $(h', d)$ is deleted.}
        \item{if $h$ is ranked lower, then the offer is ignored.}\par
    \end{itemize}

    \noindent\textbf{How many ways to match hospital-doctors?} \par
    \noindent $n!$, first doctor has $n$ choices, next doctor has $n-1$ choices, so on and so forth.
    \par

    \noindent\textbf{Will a stable match always exist?} \par
    \noindent Yes, but it is not obvious - we need to use a proof:\par
    \par
    note that during the while loop:
    \begin{itemize}
        \item{for a doctor, the hospital they go to will always be higher on their ranks.}
        \item{a hospital will be paired with doctors of decreasing ranks.}
    \end{itemize}
    If we assume that the ranking is not stable: $p = (h, d)$ and $p' = (h', d')$ such that
    $h$ prefers $d'$ and $h'$ prefers $d$.\par
    \par
    Since $h$ prefers $d'$ to $d$, then it must have made an offer to $d'$ before offering to
    $d$, but since $h$ is paired with $d$, then $d'$ would need to have rejected $h$ because
    tehy were already at a hospital they prefer to $h$. This proves that its impossible to
    have $d'$ prefer $h'$ if $h$ is ranked higher.
\end{example}

\begin{theory}{Rates of growth (Recap)} \\
    There are statements like:
    \begin{itemize}
        \item\textit{Heap sort is faster than bubble sort.}
        \item\textit{Linear search is slower than binary search.}
    \end{itemize}
    We want our language to be more precise; when is it faster, will it be not the
    case?\par\bigskip

    If we were to compare two algorithms, issues may occur!
    We mostly perfer to talk in terms of long term behaviour (asymtopics)

    For instance, if the size of an input doubles, the function could double, triple, or
    even quadruple!
\end{theory}

\begin{theory}{Big O Notation} \\
    We say that $f(n) \elemof O(g(n))$ if for large enough $n$, $f(n)$ is at most a
    constant multiple of $g(n)$ \\
    i.e. $f(n) = O(g(n))$ \\
    As such, $g(n)$ is an asymtopic upper bound for $f(n)$
\end{theory}

\end{document}
